\documentclass[12pt,a4paper]{amsart}

\usepackage{amsmath,amssymb,amsthm}
\usepackage{mathtools}
\usepackage[colorlinks=true,linkcolor=blue!60!black,citecolor=green!50!black,
            urlcolor=blue!70!black]{hyperref}
\usepackage[shortlabels]{enumitem}
\usepackage{booktabs}

%% ---- Theorem environments ----------------------------------------
\newtheorem{theorem}{Theorem}[section]
\newtheorem{lemma}[theorem]{Lemma}
\newtheorem{proposition}[theorem]{Proposition}
\newtheorem{corollary}[theorem]{Corollary}
\theoremstyle{definition}
\newtheorem{definition}[theorem]{Definition}
\newtheorem{remark}[theorem]{Remark}
\newtheorem{problem}{Problem}[section]
\newtheorem{heuristic}[theorem]{Heuristic}

%% ---- Notation ----------------------------------------------------
\newcommand{\PW}{\mathrm{PW}_\omega}
\newcommand{\Kop}{\mathcal{K}_c}
\newcommand{\GVM}{G^{(N)}_{\mathbf{p},c}}
\newcommand{\AM}{\alpha_N}
\newcommand{\BM}{\beta_N}
\newcommand{\QM}{Q^{(N)}_{\mathbf{p},c}}
\newcommand{\psin}{\psi_n^{(c)}}
\newcommand{\EN}{E_{N,c}}
\newcommand{\norm}[1]{\|#1\|}
\newcommand{\ip}[2]{\langle #1,\, #2\rangle}

\title[Coercivity and Defect Decomposition of Discrete Prolate Sampling Operators]
      {Coercivity and Defect Decomposition of Discrete Prolate Sampling Operators}

\thanks{Future work will address constructive sampling models,
sharper spectral estimates, and scaling regimes in which
both $N$ and $c$ vary.}

\author{Anonymous Author}
\address{}
\email{}
\date{\today}

\subjclass[2020]{42C05, 42B10, 47B32}

\keywords{prolate spheroidal wave functions, finite frames,
sampling in reproducing kernel Hilbert spaces,
Gram operator, Marcinkiewicz--Zygmund inequalities,
frame perturbation, coercivity}

\begin{document}
\maketitle

\begin{abstract}
Let $\PW$ be the Paley--Wiener space of $\omega$-bandlimited functions
and $V_{N,c}:=\operatorname{span}\{\psi_0^{(c)},\dots,\psi_{N-1}^{(c)}\}$
the $N$-dimensional prolate subspace at time-bandwidth product $c=\omega T$.
We study the stability of the discrete sampling form
$\QM(f)=\sum_{j=1}^N|f(p_j)|^2$ as a finite frame operator on $V_{N,c}$.

We prove that random sampling configurations and weighted sampling models
satisfying a Marcinkiewicz--Zygmund lower bound are unisolvent for $V_{N,c}$,
hence yield coercive Gram operators.
Our main structural result shows that the deviation between the discrete sampling Gram matrix
and the reference spectral matrix decomposes exactly into a weighted-frame defect
and a spectral mismatch defect.
Frame stability in the sense of a near-tight perturbation
$\QM(f)=\ip{\Kop f}{f}_{L^2}+R_N(f)$, $|R_N(f)|\le\varepsilon_N\norm{f}^2$,
follows as a corollary whenever both defects are simultaneously controlled.
All results are finite-dimensional and restricted to fixed parameters $(N,c)$.
\end{abstract}

\tableofcontents

\section{Introduction}

The prolate spheroidal wave functions (PSWFs) provide an orthonormal basis
for the finite-dimensional subspace $V_{N,c}$ of the Paley--Wiener space $\PW$
that is optimally concentrated in a time--frequency window \cite{Slepian1978,BonamiKaroui}.
From the perspective of frame theory and sampling in reproducing kernel Hilbert spaces,
the evaluation map $f\mapsto(f(p_1),\dots,f(p_N))$ defines a finite frame operator on $V_{N,c}$,
and its associated Gram matrix encodes the stability of the sampling scheme.

This paper addresses two questions.
First, which classes of sampling configurations yield coercive (invertible) Gram operators?
Second, how does the discrete sampling Gram matrix compare to the
reference spectral matrix $G^{\mathrm{ref}}_{N,c}=\Kop|_{V_{N,c}}$ in the PSWF eigenbasis?
For the first question we prove coercivity for random sampling and for weighted MZ frames.
For the second, we establish an exact algebraic decomposition of the deviation
into a frame defect and a spectral mismatch term,
which reduces the comparison problem to the simultaneous control of two explicit perturbation components.

All results are finite-dimensional and restricted to fixed parameters $(N,c)$.
The limit $N\to\infty$ and coupled limits in which both $N$ and $c$ vary are outside the scope of this paper.

\section{Setup and notation}\label{sec:setup}

\subsection{Paley--Wiener space and concentration operator}

We use the engineering Fourier convention
$\widehat{f}(\xi)=\int_{\mathbb{R}}f(x)e^{-2\pi ix\xi}\,dx$.
The \emph{Paley--Wiener space} $\PW$ consists of all
$f\in L^2(\mathbb{R})$ with Fourier support in $[-\omega/(2\pi),\omega/(2\pi)]$.
It is a reproducing-kernel Hilbert space (RKHS) with kernel
$K_{\PW}(x,y)=\tfrac{\omega}{\pi}\operatorname{sinc}(\omega(x-y)/\pi)$
and on-diagonal value $K_{\PW}(x,x)=\omega/\pi$ \cite[Eq.~(3)]{BonamiKaroui}.

\begin{remark}[Fourier convention]
Many references use $\widehat f(\xi)=\int f(x)e^{-i\xi x}\,dx$,
under which the same space has Fourier support in $[-\omega,\omega]$.
\end{remark}

The \emph{prolate concentration operator} is the self-adjoint compact operator
\[
  \Kop f(x)=\int_{-T}^{T}\frac{\sin\omega(x-y)}{\pi(x-y)}f(y)\,dy,
  \qquad \Kop:\PW\to\PW,
\]
with simple eigenvalues $1>\lambda_0(c)>\lambda_1(c)>\cdots>0$
and PSWF eigenfunctions $\{\psin\}_{n\ge 0}$ forming an orthonormal basis of $\PW$
\cite{Slepian1978,BonamiKaroui}.
We set $c=\omega T$ and
\[
  V_{N,c}:=\operatorname{span}\{\psi_0^{(c)},\dots,\psi_{N-1}^{(c)}\}.
\]
Since the PSWFs are eigenfunctions of $\Kop$, the space $V_{N,c}$ is invariant under $\Kop$
for every fixed pair $(N,c)$.

\subsection{Frame operator and Gram matrix}

\begin{definition}[Evaluation frame and Gram matrix]\label{def:GramForm}
For distinct $\mathbf{p}=(p_1,\dots,p_N)\subset[-T,T]$,
the \emph{evaluation matrix} is $M(\mathbf{p},c)_{jn}:=\psi_n^{(c)}(p_j)$.
The associated \emph{frame operator} on $V_{N,c}$ is the synthesis--analysis composition
$S_{\mathbf p}:=M(\mathbf{p},c)^*M(\mathbf{p},c)$, which we call the \emph{Gram matrix}
$\GVM:=S_{\mathbf p}$.

For $f=\sum_{n=0}^{N-1}a_n\psin\in V_{N,c}$, the identification
$f\leftrightarrow\mathbf{a}$ is unitary on $(V_{N,c},\norm{\cdot}_{L^2})$,
so $\norm{f}_{L^2}^2=\norm{\mathbf{a}}_2^2$ and
$\QM(f):=\sum_{j=1}^N|f(p_j)|^2=\ip{\GVM\mathbf{a}}{\mathbf{a}}$.
Set $\AM:=\sqrt{\lambda_{\min}(\GVM)}$, $\BM:=\sqrt{\lambda_{\max}(\GVM)}$.
\end{definition}

\begin{remark}[Spectral invariance]
The spectral quantities $\lambda_{\min}(\GVM)$ and $\lambda_{\max}(\GVM)$ are
basis-invariant on $V_{N,c}$, even though $\GVM$ itself depends on the chosen
orthonormal basis.
\end{remark}

\begin{definition}[Reference spectral matrix]\label{def:ref_matrix}
The \emph{reference spectral matrix} is the diagonal matrix
\[
  G^{\mathrm{ref}}_{N,c}:=\operatorname{diag}(\lambda_0(c),\dots,\lambda_{N-1}(c)),
\]
which is the matrix of $\Kop|_{V_{N,c}}$ in the PSWF eigenbasis.
Set $\EN:=\GVM-G^{\mathrm{ref}}_{N,c}$, the \emph{defect operator}.
\end{definition}

\begin{definition}[Frame stability condition]\label{def:T2}
We say that a sampling configuration $\mathbf{p}$ is in the
\emph{frame-stable regime} if there exists $\varepsilon_N>0$ such that
\begin{enumerate}[(S.1)]
\item $\norm{\EN}_2\le\varepsilon_N$,
\item $\varepsilon_N<\lambda_{N-1}(c)$.
\end{enumerate}
This is a near-tight frame perturbation condition;
it should not be expected to hold for generic sampling configurations.
\end{definition}

\begin{remark}[Defect operator]
Thus $\EN$ measures the deviation between the discrete sampling frame operator and
the reference spectral concentration of $\Kop|_{V_{N,c}}$.
Smallness of $\norm{\EN}_2$ is a design target, not a generic property of point sets;
one expects it, if at all, from structured quadrature or weighted MZ constructions.
\end{remark}

\subsection{A conjectural construction principle}

\begin{definition}[PSWF-adapted quadrature model]\label{def:pswf_quad}
A weighted sampling system $(\mathbf p,D)$, with $D=\operatorname{diag}(w_1,\dots,w_N)$,
is \emph{PSWF-adapted of order $N$} if
\[
\left|
\sum_{j=1}^N w_j\,\psi_m^{(c)}(p_j)\overline{\psi_n^{(c)}(p_j)}
- \lambda_n(c)\delta_{mn}
\right|
\le \varepsilon_N
\quad\text{for all }0\le m,n\le N-1.
\]
This condition asks for approximate discrete PSWF orthogonality
with respect to the eigenvalues of $\Kop$.
\end{definition}

\begin{heuristic}[Existence of PSWF-adapted systems]\label{heur:quad_candidates}
Known high-order PSWF-based quadrature constructions \cite{Xiao2001}
indicate that PSWF-adapted systems should exist for $N\lesssim c$.
However, no general existence theorem is currently available.
\end{heuristic}

\begin{remark}[Relation to the frame-stable regime]
If $(\mathbf p,D)$ is PSWF-adapted, then
$\norm{G^{\mathrm{w}}_{N,c}(\mathbf p,D)-G^{\mathrm{ref}}_{N,c}}_2 \lesssim \varepsilon_N$,
so the spectral mismatch component of the defect is controlled;
the frame-stable regime then reduces to controlling the weighted-frame defect in addition.
\end{remark}

\section{Unisolvence and frame coercivity}\label{sec:layerA}

\begin{definition}[Frame condition on $V_{N,c}$]\label{def:frame}
A sampling system $\{(p_j,w_j)\}_{j=1}^N$ is a \emph{weighted frame for $V_{N,c}$}
if there exist $0<A\le B<\infty$ such that
\[
  A\norm{f}_{L^2}^2
  \le \sum_{j=1}^N w_j |f(p_j)|^2
  \le B\norm{f}_{L^2}^2
  \qquad \text{for all }f\in V_{N,c}.
\]
When $A=B$ the frame is \emph{tight}; when $A=B=1$ it is \emph{Parseval}.
\end{definition}

\begin{definition}[Unisolvent set]\label{def:unisolvent}
A set $\mathbf{p}=(p_1,\dots,p_N)$ is \emph{unisolvent for $V_{N,c}$} if
the evaluation map $\mathcal{E}_{\mathbf{p}}:V_{N,c}\to\mathbb{C}^N$,
$f\mapsto(f(p_1),\dots,f(p_N))$, is bijective.
\end{definition}

\begin{remark}[Non-triviality of unisolvence]
For Chebyshev systems, every $N$-point set is unisolvent \cite[Ch.~1]{KarlinStudden1966}.
Since $V_{N,c}$ is not known to be a Chebyshev system in general,
unisolvence must be established separately for each sampling family.
\end{remark}

\begin{lemma}[Identity principle]\label{lem:identity}
If $f\in\PW$ vanishes on a set with a limit point in $\mathbb R$, then $f\equiv 0$.
\end{lemma}
\begin{proof}
Every $f\in\PW$ extends to an entire function of exponential type.
A non-discrete zero set forces vanishing by the identity theorem for holomorphic functions.
\end{proof}

\begin{proposition}[MZ lower bound implies unisolvence]\label{prop:mz_unisolvent}
If $\{(p_j,w_j)\}$ is a weighted frame for $V_{N,c}$ with lower bound $A>0$,
then $\mathbf p$ is unisolvent for $V_{N,c}$.
\end{proposition}
\begin{proof}
If $f(p_j)=0$ for all $j$, then $A\norm{f}_{L^2}^2\le 0$, so $f\equiv 0$ by continuity
and Lemma~\ref{lem:identity}.
\end{proof}

\begin{proposition}[Random sampling]\label{prop:random}
Let $p_1,\dots,p_N$ be i.i.d.\ with an absolutely continuous law on $[-T,T]$.
Then $\mathbf{p}$ is unisolvent for $V_{N,c}$ almost surely.
\end{proposition}
\begin{proof}
The map $(p_1,\dots,p_N)\mapsto\det M(\mathbf{p},c)$ is real-analytic.
Linear independence of $\psi_0^{(c)},\dots,\psi_{N-1}^{(c)}$ and separation of evaluation
functionals guarantee the determinant is analytically non-degenerate,
hence its zero set has Lebesgue measure zero.
Absolute continuity of the law completes the argument.
\end{proof}

\begin{lemma}[Pointwise RKHS bound]\label{lem:upper}
For $f\in\PW$ and $x\in\mathbb R$,
\[
  |f(x)|^2 \le K_{\PW}(x,x)\norm{f}_{L^2}^2 = \frac{\omega}{\pi}\norm{f}_{L^2}^2.
\]
\end{lemma}
\begin{proof}
Cauchy--Schwarz in the RKHS inner product:
$|f(x)|^2=|\ip{f}{K_x}|^2\le\norm{f}^2\norm{K_x}^2=K_{\PW}(x,x)\norm{f}^2$.
\end{proof}

\begin{theorem}[Frame coercivity]\label{thm:layerA}
Let $\mathbf{p}$ be unisolvent for $V_{N,c}$.
Then $\{f(p_j)\}_{j=1}^N$ is a frame for $V_{N,c}$ with bounds
\[
  \AM^2\norm{f}_{L^2}^2\le \QM(f)\le \frac{N\omega}{\pi}\norm{f}_{L^2}^2.
\]
In particular, coercivity holds for random sampling (a.s.) and for every MZ lower frame.
\end{theorem}
\begin{proof}
Unisolvence is equivalent to $\ker M(\mathbf p,c)=\{0\}$, hence $\GVM\succ 0$.
The lower bound follows from $\QM(f)=\ip{\GVM\mathbf a}{\mathbf a}\ge\AM^2\norm{\mathbf a}_2^2$.
The upper bound follows by summing Lemma~\ref{lem:upper} over the $N$ nodes.
\end{proof}

\begin{lemma}[Operator norm of the frame matrix]\label{lem:lambda_max_bound}
For any distinct $\mathbf p\subset[-T,T]$,
\[
  \lambda_{\max}(\GVM)
  =\sup_{\substack{f\in V_{N,c}\\ \norm{f}_{L^2}=1}} \QM(f)
  \le \frac{N\omega}{\pi}.
\]
\end{lemma}
\begin{proof}
The supremum formula is the Rayleigh quotient characterization of the largest eigenvalue
of the Hermitian matrix $\GVM$ under the PSWF isometry.
The upper bound is Theorem~\ref{thm:layerA}(upper).
\end{proof}

\begin{remark}[Sharpness]
The bound $N\omega/\pi$ is dimension-only and independent of $c$;
it is not asymptotically sharp in the PSWF regime.
Refinements using prolate concentration are left for future work.
\end{remark}

\section{Structural reduction and frame stability}\label{sec:layerB}

\begin{definition}[Weighted frame operator]\label{def:weighted_ref}
For $D=\operatorname{diag}(w_1,\dots,w_N)>0$ set
$G^{\mathrm{w}}_{N,c}(\mathbf p,D):=M(\mathbf p,c)^* D M(\mathbf p,c)$.
\end{definition}

\begin{theorem}[Defect decomposition]\label{thm:main_reduction}
For any weighted sampling model $(\mathbf p,D)$ with $D>0$,
\[
  \EN
  = M(\mathbf p,c)^*(I-D)M(\mathbf p,c)
   +\bigl(G^{\mathrm{w}}_{N,c}(\mathbf p,D)-G^{\mathrm{ref}}_{N,c}\bigr),
\]
and consequently
\[
  \norm{\EN}_2
  \le \norm{I-D}_2\cdot\norm{M(\mathbf p,c)}_2^2
   + \norm{G^{\mathrm{w}}_{N,c}(\mathbf p,D)-G^{\mathrm{ref}}_{N,c}}_2.
\]
The first term is the \emph{weighted-frame defect};
the second is the \emph{spectral mismatch defect}.
\end{theorem}
\begin{proof}
Insert and subtract $G^{\mathrm{w}}_{N,c}(\mathbf p,D)$:
\[
  \EN
  =\bigl(\GVM-G^{\mathrm{w}}_{N,c}\bigr)
   +\bigl(G^{\mathrm{w}}_{N,c}-G^{\mathrm{ref}}_{N,c}\bigr)
  = M^*(I-D)M+\bigl(G^{\mathrm{w}}_{N,c}-G^{\mathrm{ref}}_{N,c}\bigr).
\]
The norm bound follows from subadditivity and submultiplicativity of $\norm{\cdot}_2$.
\end{proof}

\begin{remark}[Algebraic vs.\ probabilistic coupling]
Both defect terms share the same evaluation matrix $M(\mathbf p,c)$.
The decomposition is algebraic, not probabilistically or geometrically decoupled;
controlling one term does not automatically control the other.
\end{remark}

\begin{corollary}[Frame-stable regime from defect bounds]\label{cor:two_defects}
Suppose
\[
  \norm{I-D}_2\le \delta_N,
  \qquad
  \norm{G^{\mathrm{w}}_{N,c}(\mathbf p,D)-G^{\mathrm{ref}}_{N,c}}_2\le \eta_N,
  \qquad
  \varepsilon_N:=\delta_N\tfrac{N\omega}{\pi}+\eta_N < \lambda_{N-1}(c).
\]
Then $\mathbf p$ is in the frame-stable regime with parameter $\varepsilon_N$.
\end{corollary}
\begin{proof}
By Theorem~\ref{thm:main_reduction} and Lemma~\ref{lem:lambda_max_bound},
$\norm{\EN}_2\le\delta_N\cdot N\omega/\pi+\eta_N=\varepsilon_N<\lambda_{N-1}(c)$.
\end{proof}

\begin{lemma}[Weyl perturbation bound]\label{lem:weyl_stability}
In the frame-stable regime,
\[
  \lambda_{\min}(\GVM)\ge\lambda_{N-1}(c)-\varepsilon_N>0.
\]
\end{lemma}
\begin{proof}
Apply the Hermitian perturbation inequality
$\lambda_{\min}(A+E)\ge\lambda_{\min}(A)-\norm{E}_2$
with $A=G^{\mathrm{ref}}_{N,c}$ and $E=\EN$.
\end{proof}

\begin{corollary}[Frame stability]\label{cor:spectral_comparison}
In the frame-stable regime,
\[
  \QM(f)=\ip{\Kop f}{f}_{L^2}+R_N(f),
  \qquad |R_N(f)|\le\varepsilon_N\norm{f}_{L^2}^2.
\]
No global analytic continuation or spectral equivalence beyond finite-dimensional
perturbation theory is claimed.
\end{corollary}
\begin{proof}
In the PSWF eigenbasis,
$\ip{\Kop f}{f}_{L^2}=\ip{G^{\mathrm{ref}}_{N,c}\mathbf a}{\mathbf a}$.
So
$\QM(f)=\ip{G^{\mathrm{ref}}_{N,c}\mathbf a}{\mathbf a}+\ip{\EN\mathbf a}{\mathbf a}$.
Setting $R_N(f):=\ip{\EN\mathbf a}{\mathbf a}$ gives
$|R_N(f)|\le\norm{\EN}_2\norm{\mathbf a}_2^2\le\varepsilon_N\norm{f}_{L^2}^2$.
\end{proof}

\section{Open problems}\label{sec:open}

\begin{problem}[Explicit PSWF frames]
Construct sampling sets and positive weights satisfying a Marcinkiewicz--Zygmund frame condition on $V_{N,c}$
with explicit bounds $(A,B)$.
\end{problem}

\begin{problem}[Uniform-grid unisolvence]
Determine whether uniform grids are unisolvent for $V_{N,c}$.
\end{problem}

\begin{problem}[Optimal weighting for near-diagonalization]
Derive quantitative estimates on the weighted-frame defect and spectral mismatch defect
sufficient to trigger Corollary~\ref{cor:two_defects}.
\end{problem}

\begin{problem}[Coupled limits]
Analyze $\lambda_{\min}(\GVM)$ and the competition between
$\norm{\EN}_2$ and $\lambda_{N-1}(c)$ in regimes where $N$ and $c$ vary simultaneously.
\end{problem}

\begin{thebibliography}{99}

\bibitem{BonamiKaroui}
A.~Bonami and A.~Karoui,
\textit{Spectral decay of the sinc kernel operator and approximation
with prolate spheroidal wave functions},
arXiv:1012.3881.

\bibitem{HornJohnson2013}
R.~Horn and C.~Johnson,
\textit{Matrix Analysis}, 2nd ed.,
Cambridge University Press, 2013.

\bibitem{KarlinStudden1966}
S.~Karlin and W.~Studden,
\textit{Tchebycheff Systems},
Interscience, 1966.

\bibitem{Slepian1978}
D.~Slepian,
\textit{Prolate spheroidal wave functions, Fourier analysis,
and uncertainty~-- V},
Bell Syst.\ Tech.\ J.\ \textbf{57} (1978), 1371--1430.

\bibitem{Xiao2001}
H.~Xiao, V.~Rokhlin, and N.~Yarvin,
\textit{Prolate spheroidal wave functions, quadrature, and interpolation},
Inverse Problems \textbf{17} (2001), 805--838.

\end{thebibliography}

\end{document}
